\chapter{System Overview and Problem Formulation\label{cha:chapter4}}


This chapter provides the system-level context and problem formulation for the theoretical developments of Chapter 5 (steady-state design) and Chapter 6 (fault-tolerance design). Section~4.1 describes the warehouse logistics task, and formulates the multi-agent coverage problem. Section~4.2 specifies the capabilities of the robotic fleet at the agent level. Section~4.3 discusses the UAV-proxy paradigm enabling vertical scanning of tall shelves. Section~4.4 briefly describes the simulation environment in which the system was experimented and validated, while deferring the implementation-level details to Chapter 7.

\section{Global Mission Context: Warehouse Logistics and Corridor Constraints}

Table~\ref{tab:lit_review_comparison} concluded the preceding chapter with ten specific deltas between the surveyed literature and the system under investigation; the current one addresses each of these discrepancies as a particular design commitment. Given Chakraa et al.\ and Quinton et al.'s findings on auction-based and optimization-based task allocation mechanisms' limited scalability and communication assumptions \cite{chakraa2023optimization,quinton2023market} and Agrawal et al.'s DC-MRTA's increased computational complexity due to its coupling with traffic-aware navigation \cite{agrawal2022dcmrta}, the disjoint task partition $\mathcal{T}_i(t) \subset \mathcal{P}$ formalized later in this section (Eq.~\eqref{eq:ch4_coverage_partition}) is instead addressed by a static, near-costless cardinal sweep, as described in Section~5.3. Similarly, Sharma et al.'s dynamically re-partitioned DARP regions \cite{sharma2026darp} are replaced, for the fleet $\mathcal{A}$ (Eq.~\eqref{eq:ch4_fleet}), with a fixed \texttt{owned\_regions} assignment established at deployment (Section~5.3.1), the losses of parallelism for which will be empirically quantified in Chapter~8. The narrow-aisle bottleneck formally characterized later in this section is another area of reduced complexity: the formally verified reservation graphs and token circulation mechanisms of Grover et al., Pratissoli et al., He et al., and Parihar and Chakraborty \cite{grover2023deadlock,pratissoli2024hierarchical,he2023deadlock,parihar2022token,lee2025optimizing} are instead addressed by a simplified, unverified TTL-leasing and LCFS-overwriting peer broadcasting mechanism described in Section~5.2. Fault detection adheres to a similar paradigm of reduced autonomy. Lee et al.\ and Bossens et al.\ both recommend continuously tracking the behavior model of a robot navigating its task \cite{lee2022datadriven,bossens2022resilient}; the fail-stop threshold $E_{\text{fail}} = 0.0\,\%$ introduced in Section~4.2.4 below is instead watched over by two discrete channels: the self-declaration watchdog and the heartbeat timeout mechanism of Section~6.1, neither of which attempts behavioral inference on a robot unless contact has already been lost. Reallocation similarly favors simplicity: the continuously reoptimized reward vectors of Mayya et al.'s MRTA and Gong et al.'s clustered, locally updated multi-robot planning \cite{mayya2021resilient,gong2024resilient} are replaced, for the single-fleet setting, with a deterministic, battery-proportional split (Section~6.6), triggered strictly per confirmed failure instead of being recomputed on a rolling schedule. Localization recovery narrows Shi et al.'s visual auxiliary positioning \cite{shi2022vision} and the particle filter maintenance procedure reviewed by Macenski et al.\ \cite{macenski2023fromdesks} to the covariance-gated two-stage spinning and resetting procedure of Section~6.2, which waits for divergence rather than forestalling it. The same localization failure response channel enables the same fail-stop commitment to be exercised across the fleet to address the energy-aware allocation gap: Notomista et al.\ and Djenadi and Mendil\ both recommend optimizing charge-preserving trajectories assuming the fleet robots have a return-to-dock opportunity \cite{notomista2022resilient,djenadi2022energyaware}, which Section~4.2.4 explicitly rejects---leaving Sections~6.1.1 and~6.6 to handle total battery exhaustion as an expected outcome rather than a condition to navigate around. The perception stack is the only aspect of this system that differs from the comparison literature due to practical rather than conceptual reasons. Contrary to the CNN and transformer-based defect classifiers, accelerator-dependent processing pipelines, and 3D scans from physical UAVs reviewed by Hussain and Hill, Sen et al., Benmachiche et al., and Tsiakas et al.\ \cite{hussain2023custom,sen2025feature,benmachiche2026realtime,tsiakas2024autonomous}, the handcrafted LBP/HOG-plus-linear-SVM pipeline reviewed in Section~5.5 is constrained to the CPU resources of the WSL2 host, a design choice that will be further qualified in this chapter through the sensor and UAV-specific constraints outlined in Sections~4.2.3 and~4.3. Explainability is a direct result of the perception stack's reduced autonomy: Section~6.7's closed-form back-projection and untested SVM/XGBoost fallback step around the model-agnostic post-hoc explanation techniques reviewed by Sakai and Nagai\ and Myllyaho et al.\ \cite{sakai2022explainable,myllyaho2022misbehaviour}, offering lower runtime costs at the expense of generality---a trade-off Chapter~8 re-examines directly. Finally, the general-purpose ROS~2 benchmarking platforms and default path tracking controllers reviewed by Macenski et al.\ and Shcherbyna et al.\ \cite{macenski2023fromdesks,shcherbyna2024arena,macenski2023regulated} are replaced, in Section~7.2 and the simulation substrate described in Section~4.4 below, with a custom, gate-based readiness polling system; this chapter's contribution was to identify the specific requirements a three-robot fleet simulation had to impose on a shared host rather than build upon the existing research aiming to generalize across arbitrary multi-robot settings. None of these ten design choices are presented as improvements over the state-of-the-art methods they are motivated by; rather, they are formulated as implementation realities to be acknowledged before moving onto evaluating the empirical performance of the system under the constraints defined by it, the subject of Chapter~8.

A critical task of any cyber-physical logistics facility is the inventory auditing procedure, which requires a systematic examination of the stored items. This process usually involves either the manual counting of the objects by human workers or the automated analysis of a visual scan of a large-scale inventory storage area. Both methods involve substantial safety, time cost, and human labor expenditures, especially when examining targets located on large high-rise shelves. Autonomous robotic platforms offer tremendous promise for such logistics tasks, but the use of decentralized multi-robot coalitions introduces substantial additional spatial coordination challenges.

The physical warehouse environment used as the basis for this research was designed to resemble a realistic logistics facility. The workspace consists of several structural pillars, dynamic loading docks, and multiple rows of storage shelves and pallets, which were placed to host inventory targets. In particular, to assign the robots global navigation waypoints, the warehouse workspace was represented as a bounded 2D space, as defined in Equation~\eqref{eq:ch4_workspace}:

\begin{equation}
\mathcal{W} = [X_{\min}, X_{\max}] \times [Y_{\min}, Y_{\max}] \subset \mathbb{R}^2 \label{eq:ch4_workspace}
\end{equation}

where $\mathcal{O}$ is the set of obstacles in Equation~\eqref{eq:ch4_obstacles} defining the free workspace for robot motions, where elements $\mathcal{O}_i$ are members of the power set $\mathcal{P}(\mathcal{W})$:

\begin{equation}
\mathcal{O} = \{ \mathcal{O}_1, \mathcal{O}_2, \dots, \mathcal{O}_M \} \subset \mathcal{P}(\mathcal{W}) \label{eq:ch4_obstacles}
\end{equation}

The resulting free workspace $\mathcal{W}_{\text{free}}$ available for navigation is defined according to Equation~\eqref{eq:ch4_free_workspace}:

\begin{equation}
\mathcal{W}_{\text{free}} = \mathcal{W} \setminus \bigcup_{i=1}^{M} \mathcal{O}_i \label{eq:ch4_free_workspace}
\end{equation}

\subsection{Physical Workspace Description}

The workspace consists of twenty-two large, small shelves, and pallets arranged in a manner to represent high-density storage units. Figure~\ref{fig:ch4_warehouse_overview} contrasts this arrangement as rendered in the Gazebo Harmonic client, showing both the plan-view layout of the storage rows and the oblique perspective a TurtleBot3 Waffle unit encounters while traversing the central aisle.

\begin{figure}[htbp]
\centering
\subcaptionbox{Plan-view of the warehouse layout, showing the twenty-two shelf and pallet units alongside the structural pillars.\label{fig:ch4_warehouse_top}}
    {\includegraphics[width=0.47\textwidth]{fig_4_1_1_gazebo_world_top_view.pdf}}
\hfill
\subcaptionbox{Oblique ground-level perspective of the central aisle, corresponding to the viewpoint of an active agent.\label{fig:ch4_warehouse_back}}
    {\includegraphics[width=0.47\textwidth]{fig_4_1_1_gazebo_world_back_view.pdf}}
\caption{The simulated warehouse environment rendered in Gazebo Harmonic.}
\label{fig:ch4_warehouse_overview}
\end{figure}

\subsubsection{Discretization of Storage Units into Inventory Cells}

Each of the storage unit types contains several inventory cells located at various heights along the length of the facility:

\begin{equation}
\mathcal{P} = \{ p_1, p_2, \dots, p_{142} \}
\end{equation}

Each target location $p_i$ therefore represents a discrete inspection point in the free space, defined as:

\begin{equation}
p_i = (x_i, y_i, z_i, \theta_i) \in \mathcal{W}_{\text{free}} \times \mathbb{R} \times S^1
\end{equation}

where the first two coordinates define the position, while the last one encodes the yaw angle orthogonal to the surface of the cardboard box, allowing the agent to capture high-resolution imagery of the target surface. The coordinates of the target locations were hardcoded into the world geometry file \texttt{tugbot\_warehouse\_clean.sdf}.

\subsubsection{Parsing of the World Geometry}

To generalize the multi-agent planning procedure to arbitrary workspace layouts without hardcoding the target locations, the planning algorithm parses the world geometry \texttt{sdf\_parser.py}. In particular, upon initializing the simulation world from the XML description \texttt{tugbot\_warehouse\_clean.sdf}, the parser extracts the geometric centroids of the storage units and calculates approximate approach waypoints for the robotic agents. This procedure therefore allows the planner to adapt to various warehouse layouts.

\subsection{Aisle and Corridor Constraints}

The storage shelves are organized in parallel rows, forming navigation corridors:

\begin{equation}
\mathcal{C} = \{ c_1, c_2, \dots, c_K \}
\end{equation}

which contain the majority of the workspace obstacles.

\subsubsection{Narrow Aisle Bottlenecks}

The width of the corridor $W_{\text{aisle}}$ is insufficient to allow the autonomous agents to pass in both directions:

\begin{equation}
W_{\text{aisle}} < 2 \cdot D_{\text{safe}} + \epsilon
\end{equation}

where $D_{\text{safe}}$ encodes the clearance distance required around the robot to prevent collisions with other agents. In particular, if two agents attempt to move from opposite ends of the narrow aisle, their localized costmaps will inflate the obstacle boundaries, preventing mutual passage:

\begin{equation}
D_{\text{safe}} + D_{\text{robot}} > W_{\text{aisle}}
\end{equation}

rendering the narrow aisle a global coordination bottleneck for the multi-robot system \cite{matos2025efficient}. 

\begin{figure}[h]
    \centering 
    \includegraphics[width=0.7\textwidth]{img/fig_4_1_corridor_bottleneck.pdf} 
    \caption{Geometric configuration of a narrow aisle bottleneck, showing $W_{\text{aisle}}$ relative to $2 \cdot D_{\text{safe}}$.} 
    \label{fig:fig_4_1_corridor_bottleneck} 
\end{figure}
Figure~\ref{fig:fig_4_1_corridor_bottleneck} shows this effect for a generic narrow aisle navigation corridor. Figure~\ref{fig:ch4_narrow_aisle_physical} shows the same bottleneck as it occurs inside the simulated environment, with the robot occupying the greater part of the available lateral clearance while approaching a small-shelf inspection face.

\begin{figure}[htbp]
\centering
\subcaptionbox{Approach view of a TurtleBot3 Waffle unit entering a narrow aisle between two facing small-shelf units.\label{fig:ch4_narrow_aisle_a}}
    {\includegraphics[width=0.47\textwidth]{fig_4_1_2_narrow_aisle_1.pdf}}
\hfill
\subcaptionbox{The same aisle viewed from a closer vantage point, illustrating the residual clearance margin either side of the chassis.\label{fig:ch4_narrow_aisle_b}}
    {\includegraphics[width=0.47\textwidth]{fig_4_1_2_narrow_aisle_2.pdf}}
\caption{Physical realization of the narrow aisle bottleneck of Figure~\ref{fig:fig_4_1_corridor_bottleneck} inside the Gazebo simulation.}
\label{fig:ch4_narrow_aisle_physical}
\end{figure}

\subsubsection{Interference With Building Pillars}

The facility contains several steel building pillars, located at approximately $Y \approx -15.00\text{ m}$ that interfere with the default approach waypoints for the large shelves. In particular, the naive approach of placing the inspection waypoints at a constant distance from the side of the large storage unit places the agent's trajectory inside the inflation zone of the steel pillars. Therefore, an additional procedure must be implemented to allow the robotic agents to bypass the building pillars during the inspection of the large storage units.

\subsection{Inspection Target Distribution and Multi-Agent Coverage}

The global mission requires the agents to audit all $|\mathcal{P}| = 142$ inspection targets distributed across the storage units. In particular, the multi-robot audit procedure should ensure complete target set coverage while minimizing the total mission duration.

The total target set, $ |\mathcal{P}| = 142 $, is developed by enumerating the number of inspection cells in all small static shelving units (8 cells / shelf, utilizing the \texttt{yminus} and \texttt{yplus} approach sides), large static shelves (8 cells / shelf, subject to unidirectional \texttt{xminus} accessibility limits enforced on \texttt{shelf\_big\_1}) and static pallets (1 target / pallet). The dynamic mobile pallets provide 6 additional targets by including the two-sided (\texttt{xminus}/\texttt{xplus}) 3-level section of the integrated \texttt{mobile\_cluster}.


\subsubsection{Multi-Agent Fleet Description}

The robotic fleet consists of three homogeneous differential drive agents, formalized as the set $\mathcal{A}$ in Equation~\eqref{eq:ch4_fleet}:

\begin{equation}
\mathcal{A} = \{ a_1, a_2, a_3 \} \label{eq:ch4_fleet}
\end{equation}

which operate in a fully decentralized manner. At any given moment $t$, the state of each active agent $a_i \in \mathcal{A}$ is defined by its pose in the free space $x_i(t) \in \mathcal{W}_{\text{free}}$. Collision avoidance between active physical frames is mathematically governed by the clearance constraint in Equation~\eqref{eq:ch4_collision_clearance}:

\begin{equation}
\text{For all } a_i, a_j \in \mathcal{A}, a_i \neq a_j: \quad \|x_i(t) - x_j(t)\| \geq 2 \cdot r_{\text{bounding}} \label{eq:ch4_collision_clearance}
\end{equation}

where $r_{\text{bounding}}$ represents the conservative bounding radius circumscribing the rectangular physical footprint of the chassis.

\subsubsection{Target Set Coverage Task Formulation}

Let $\mathcal{T}_i(t) \subset \mathcal{P}$ denote the set of targets assigned to the robotic agent $a_i$ at time $t$. In particular, the multi-agent task allocation procedure should satisfy the coverage and disjointness constraints in Equation~\eqref{eq:ch4_coverage_partition}~\cite{xu2026energy}:

\begin{equation}
\forall t: \bigcup_{i=1}^{3} \mathcal{T}_i(t) = \mathcal{P} \quad \text{and} \quad \mathcal{T}_i(t) \cap \mathcal{T}_j(t) = \emptyset \quad \forall i \neq j \label{eq:ch4_coverage_partition}
\end{equation}

The mission is considered complete when all targets have been audited. In particular, the global mission completion time $T_{\text{audit}}$ should be minimized across all valid task allocations, representing the time when the last agent finishes its task queue and returns to its docking pose.

The main challenge of the decentralized multi-agent coverage task is the absence of the master planner, which would coordinate the agents and resolve the spatial coordination conflicts described in Section~4.1.2. In particular, the agents should negotiate the access to the narrow aisle using the ad-hoc communication described in \texttt{waypoint\_sender.py}, resolving the spatial inter-agent assignment conflicts while following the constraints imposed by the inspection waypoints. Furthermore, in the case of a single agent failure, the surviving agents should be able to reconfigure the task allocation table to reassign the inspection targets of the failed robot \cite{gong2024resilient}.

\subsubsection{Mobile Clustering for Dynamic Inventory Asset Grouping}\label{sec:ch4_mobile_cluster}

Dimensional configuration cost incurred by multi-agent task allocation is avoided by dynamically grouping the non-stationary inventory assets located in the logistics facility. Namely, particular mobile pallet entities, identified as \texttt{pallet\_box\_mobile}, \texttt{pallet\_box\-\_mobile\_0}, and \texttt{pallet\_box\_mobile\_1}, display extremely unstable positions in relation to the logistics process. To properly group these dynamic agents for the decentralized sweep protocol, the spatial parser utilizes an artificial clustering technique. The centroids of the detected mobile pallets are grouped to form one virtual object called \texttt{mobile\_cluster}.

The coordinates of this object are represented through the geometric mean of the centroids of the corresponding pallets:

\begin{equation}
c_{\text{cluster}} = \left( \frac{1}{M} \sum_{k=1}^{M} x_k, \quad \frac{1}{M} \sum_{k=1}^{M} y_k \right)
\end{equation}

where $M = 3$ denotes the number of the mobile pallets in the subset. Using the generation of unified orthogonal poses of the cluster segments along the global $X$-axis and a deterministic North-to-South sweeping order (by descending $Y$-coordinates), the multi-agent task allocator avoids local planning issues and redundant entries of the sweep corridors.



\section{Homogeneous Fleet Description and Agent Capabilities}

To reduce the complexity of multi-agent coordination while maximizing the reliability and flexibility of the robotic fleet, this research uses a homogeneous fleet of three differential drive robots. In particular, the failure of any single robot can be compensated by the other two, while the uniformity of the dynamic and sensory properties allows the agents to perform any inspection tasks interchangeably without additional calibration.

The robotic agents were modeled after the TurtleBot3 Waffle differential drive mobile robot.

\subsection{Mechanical and Kinematic Properties}

The differential drive robots have standard 2D TurtleBot3 Waffle specifications, with the track width and wheel radius defining the robot's kinematic constraints. In particular, these configuration parameters are loaded dynamically at runtime from the agent's parameter configuration space.

\subsubsection{Track Width (Wheel Separation)}

The lateral distance between the left and right wheels is configured as:

\begin{equation}
w_{\text{track}} = 0.287\text{ m}
\end{equation}

The distance between the robot's left and right wheels defines its turning capabilities: a larger value of $w_{\text{track}}$ increases the minimal turning radius of the robot.

\subsubsection{Wheel Radius}

The radius of the drive wheels is:

\begin{equation}
r_{\text{wheel}} = 0.033\text{ m}
\end{equation}

The wheel radius determines the ratio between the angular velocity of the robot's motors and its linear speed.

\subsection{Dynamic Constraints}

To avoid mechanical damage, the execution engine limits the linear velocity and acceleration of the robots:

\begin{equation}
v_{\text{max}} = 0.8\text{ m/s}
\end{equation}

\begin{equation}
a_{\text{max}} = 1.0\text{ m/s}^2
\end{equation}

The odometry update loop runs at $50\text{ Hz}$, providing sufficient control authority to the low-level navigation stack.

\subsection{Sensor Suite and Spatial Perception}

Each robot was equipped with the same sensor suite to ensure reliable localization and inventory inspection capabilities.

\subsubsection{Extended-Range Laser Scanner}

In standard ROS 2 Gazebo simulations, 2D laser rangefinders have a limited $3.5\text{ m}$ range, which creates localization challenges in large-scale facility layouts. In particular, because the robot cannot observe the opposite walls of the central courtyards, the AMCL localization filter produces large pose estimation errors.

To address this issue, the warehouse workspace file overrides the default maximum range of the onboard 2D laser scanner to $20.0\text{ m}$. In particular, the sensor processes $360$ angular bins, providing $1^\circ$ angular resolution over an extended $20.0\text{ m}$ range. This allows the robot localization filter to observe the building pillars at the center of the simulation world, reducing the drift of the AMCL particle filter \cite{wu2025analysis}.

\subsubsection{Optical Sensor and Vertical Scanning}

Each robot has a forward-looking optical sensor, capable of capturing $128 \times 128$ grayscale images of the storage unit surface. To allow the agents to inspect the tall storage units, the robot control algorithm uses a time-based state machine to emulate UAV altitude control. In particular, while approaching the large storage unit, the robot interrupts the global navigation planner and enters the vertical scanning state. This state artificially increases the robot's Z-position (by adding a $4.0\text{ s}$ time delay for big shelves and $1.5\text{ s}$ delay for small shelves), allowing it to capture an image of the target surface with the onboard optical sensor following the geometric parameters parsed dynamically at runtime as formulated in Section~5.3.

\subsection{Battery Limits and Hybrid Power Model}

The autonomy of each robotic agent is defined by its internal battery capacity. In particular, the fleet's power management module monitors the State of Charge (SOC) of each robot:

\begin{equation}
E_{\text{total}}(t) \in [0.0\%, \; 100.0\%]
\end{equation}

\subsubsection{Hybrid Discharge Model}

Most autonomous mobile robots use a linear time-based discharge model, where SOC decreases with a constant rate after reaching the maximum capacity. In practice, however, the discharge process follows a more complex hybrid discharge pattern, governed by static baseline idling drain and dynamic, motion-induced load\cite{liu2023open}:

\begin{equation}
E_{\text{total}}(t) = E_{\text{start}} - \left( E_{\text{idle}}(t) + E_{\text{motion}}(t) \right)
\end{equation}

The static component approximates the power consumed by the computer vision and perception modules. The dynamic component is proportional to the linear velocity of the robot:

\begin{equation}
E_{\text{motion}}(t) \propto \int_{0}^{t} v(\tau) \, d\tau
\end{equation}

Therefore, robots that move at higher speeds (e.g., to avoid the narrow aisle bottlenecks) will experience faster battery drainage during the inspection of the inventory targets. This effect becomes most prominent during the long-range navigation tasks required to reach the distant inspection waypoints.

\subsubsection{Low-Power Threshold and Fail-Stop Depletion Paradigm}

In order to assess the properties of the decentralized re-allocation and dynamic recovery procedures for the fleet under the worst-case operating scenarios, the return-to-dock procedure of the power management module is entirely disengaged. Autonomous low-power return navigation attempts in such scenarios pose a threat to the operational integrity of the entire fleet due to high likelihood of mid-transit battery depletion in narrow-aisle multi-agent warehouse environments.

As such, the physical simulation environment implements a fail-stop depletion paradigm. The state of charge is tracked until:

\begin{equation}
E_{\text{fail}} = 0.0\% \label{eq:absolute_failure_threshold}
\end{equation}
in which case the agent is forced to perform an immediate physical stop and transition to the \texttt{FAILED} state. The abrupt termination is necessary to evaluate the individual agent's capacity for fault detection and subsequent spatial claim release and proportional split re-allocation procedures described in detail in Chapter~\ref{cha:chapter6}.

\subsection{Edge Compute Architecture}

The decentralized nature of the multi-robot system imposes additional constraints on the compute architecture. In particular, each robot operates as an edge compute node, locally processing the sensor data and maintaining its own localization estimate. In particular, the namespace isolation procedure and the underlying ROS 2 peer-to-peer communication protocol described in Section~5.2 and Section~7.2 enable the edge-deployed autonomy stack to satisfy the decentralized coordination requirements.

\section{UAV Proxy Paradigm: Justification for Z-Axis Simulation}\label{sec:ch4_uav_proxy}

Before proceeding to the argument below, the qualification to the following claim requires stating: the author is not asserting that the mechanism described in this work equips the TurtleBot3 Waffle platform with rotorcraft propulsion or any other form of physical flight capability, since similar facilities in warehouse settings perform orthogonal high-shelf viewing by employing a scissor-lift module, a fixed mast-mounted sensor column, or a pan-tilt camera head instead of free-flight capabilities \cite{tsiakas2024autonomous}, any of which could be used as a substitute for the mechanism under consideration. What follows is a heuristic that entirely abstracts away from the choice of mechanism, representing a time spent by a vertical-access module in repositioning as a fixed dwell superimposed on top of a planar $\text{SE}(2)$ motion as described in the original Chapter~\ref{cha:chapter2}, so that the contribution of the present work to the coordination, allocation, and fault recovery can be assessed independent from that choice. Section 4.3.3 will return to this qualification once the mathematical foundation for orthogonal viewing has been laid.

Inspecting the cardboard surface of the storage unit requires the robotic agent to capture high-resolution images from a sufficient distance away while maintaining orthogonal vision on the target plane. In an industrial logistics warehouse, the storage units are placed on large-scale $1.8\text{ m} \times 3.0\text{ m} \times 6.0\text{ m}$ shelves, which extend vertically along the Z-axis:

\begin{equation}
Z_{\text{shelf}} \in [1.8\text{ m}, \; 6.0\text{ m}]
\end{equation}

Conventional differential drive robots are limited to planar 2D navigation and cannot change their Z-position during the inspection of the storage shelf.

To simplify the control algorithm while retaining the ability to capture images of the vertical storage shelves from an orthogonal viewpoint, this project simulates a UAV proxy paradigm for the multi-robot system.

\subsection{Mathematical Justification for UAV Proxy}

The ability to capture high-resolution images of the storage unit surface is of particular importance for the anomaly detection component of this project. In particular, the local binary patterns descriptor is highly sensitive to perspective distortions caused by non-orthogonal projection of the target surface on the image plane.

Let us define the vertical storage shelf as a flat surface aligned along the Z-axis at a distance $D$ from the robot's optical sensor. For generality, let us assume that the target surface is located at an arbitrary height $H_{\text{cell}}$ relative to the robot's center of mass, and that the robot optical sensor is placed at a height $h_{\text{camera}}$.

The robot attempting to capture the image of a high-rise cell must therefore observe the target surface at an angle $\theta_{\text{camera}}$ to the horizontal plane:

\begin{equation}
\theta_{\text{camera}} = \arctan\left(\frac{D}{H_{\text{cell}} - h_{\text{camera}}}\right)
\end{equation}

which would cause significant perspective warping of the target image:

\begin{equation}
\theta_{\text{camera}} \neq 90^\circ \quad \forall H_{\text{cell}} \neq h_{\text{camera}}
\end{equation}

Therefore, if the robot is not capable of changing its vertical position, the Z-axis location of the target surface determines the skew of the perspective projection observed by the optical sensor. This effect would be particularly detrimental for the edge detection component of the vision-based anomaly detection algorithm, potentially causing false positives due to perspective-induced texture distortion, as shown in [figure2]

\begin{figure}[h]\centering \includegraphics[width=0.95\textwidth]{img/fig_4_2_camera_angle_geometry.pdf} \caption{Incidence angle $\theta_{\text{camera}}$ as a function of camera height $h_{\text{camera}}$ and target cell height $H_{\text{cell}}$; convergence to $90^\circ$ under the UAV-proxy vertical translation.} \label{fig:camera_angle} \end{figure} 

To remove this source of perspective distortion, the robotic agent must be able to observe the target surface at a viewing angle of $\theta_{\text{camera}} = 90^\circ$.

Under a UAV proxy paradigm, we can assume that the robotic agent is capable of freely adjusting its Z-position to meet this requirement:

$h_{\text{camera}} \to H_{\text{cell}}$

Substituting this value into the expression for $\tan(\theta_{\text{camera}})$, we obtain:

\begin{equation}
\theta_{\text{camera}} = \arctan\left(\frac{D}{0}\right)
\end{equation}

\begin{equation}
\theta_{\text{camera}} = 90^\circ
\end{equation}

This perspective projection eliminates any skew, enabling the image resolution algorithm to correctly identify the micro-texture patterns of the cardboard surface.

\subsection{Mechanical Proxy and State Machine}

Instead of implementing a full 3D flight dynamics model, which would require defining aerodynamic parameters, rotor torque characteristics, and wind effects, we can simply emulate the UAV's vertical maneuver using a state machine.

In particular, while approaching the target inspection waypoint, the robot interrupts its global navigation planner and advances into the vertical scanning state. This state adds a time-dependent delay $\Delta t$ to the forward motion of the robot, which approximates the increase in UAV altitude during the vertical climb:

\begin{equation}
\Delta t = \begin{cases} 4.0\text{ s} & \text{for large shelves} \\ 1.5\text{ s} & \text{for small shelves} \end{cases}
\end{equation}

During this time interval, the robot halts its horizontal motion (by setting the angular and linear velocities to zero) and captures a $128 \times 128$-pixel image of the target surface with the onboard optical sensor. The delay duration $\Delta t$ is chosen to be proportional to the height of the storage shelf parsed from the world SDF file, allowing the robot to emulate a continuous increase in the Z-position during the inspection of multi-tier shelves.

\subsection{Limitations of the UAV Proxy Paradigm}

While the described approach indeed allows the robotic agent to capture orthogonal images of the storage shelf without incurring additional mechanical complexity, it does not account for a few practical effects:

\begin{itemize}
    \item Aerodynamic ground effects and wind created by the rotation of the rotor blades,
    \item Additional battery drain from the increased power draw by the UAV motors,
    \item The physical embodiment of the described vertical reach: a scissor-lift mechanism, a fixed-height sensor mast, or a pan-tilt camera on a fixed-height mast would have different degrees of mechanical complexity, space requirements, and control authority, which are outside the scope of this thesis (as stated in the beginning of this section and in Section~\ref{sec:ch1_scope}).
\end{itemize}

The paradigm described in this study can thus be seen as an abstraction that hides implementation-specific details from the robotic autonomy stack, be it a scissor-lift mechanism or a camera gimbal. The contributions of this thesis to the autonomy stack described in Chapters~\ref{cha:chapter5} through~\ref{cha:chapter8} are independent of the implementation and hold regardless of the choice of mechanism.

\section{Simulation Environment Description}

The simulation experiments were carried out in a GPU-accelerated environment running on Windows Subsystem for Linux 2 (WSL2), which hosted Ubuntu 24.04 LTS. Specifically, ROS 2 Jazzy Jalisco and the Gazebo Harmonic simulation software \cite{afzal2021gzscenic} with the Ogre2 rendering engine were utilized for physics simulation acceleration and high-quality computer graphics, respectively. The GPU acceleration was necessary to attain the Real-Time Factor of $\text{RTF} \approx 1.0$. The specifics of graphics acceleration, window manager, and UI rendering fall outside the scope of this research and are detailed in Section~7.1 and Section~7.4. In particular, the timing anomalies introduced by the WSL2 hypervisor are relevant to the fault detection threshold, which is detailed in Section~6.1.